\section{Conclusion}
\label{sec:conclusion}

Our central finding is that return, ranking accuracy, and regulatory risk suitability need not conflict on FAR-Trans: a recommender can be made far more suitable to the customer's declared risk band while keeping, and even improving, accuracy and return. This makes suitability an architectural choice rather than a tradeoff patched by post-hoc re-ranking.

The results assemble a diagnostic-to-remedy chain. An analytical study establishes that discordance is prevalent and a persistent customer-level trait, and that coherent buys earn a statistically significant return premium, removing the economic objection to optimising for suitability. Auditing the FAR-Trans baselines on the three axes shows that both under-serve specific risk bands and that neither balances the three. Risk-based segmentation then closes the per-band deficit. Of its two exclusive levers, regrouping customers onto their revealed band is the dominant one: it lifts the harmonic-mean balance above the baseline by raising suitability and accuracy together at no cost to return. The margin loss, by contrast, buys suitability by spending accuracy, and so lowers the balance.

More broadly, the approach shows that regulatory suitability constraints can be encoded as training-time objectives and, better still, corrected at the data source rather than filtered afterwards which opens a path for recommenders in other regulated domains, such as insurance product suitability.
